\section*{Anexo 4-G --- Catálogo de interfaces internas}
\phantomsection
\label{anx:G}
\addcontentsline{toc}{section}{Anexo 4-G --- Catálogo de interfaces internas}
El catálogo conserva los identificadores de las ocho interfaces internas. Las fichas que siguen completan el modo, volumen, ventana y conducta ante lentitud exigidos por RT-05.21. Los volúmenes son escenarios de diseño, no mediciones. Un mensaje de negocio puede atravesar varias interfaces: sus cifras no se suman como operaciones únicas.

\begin{tablalafrox}{Integraciones internas: dueño, contrato y continuidad}{tab:catalogo-internas}{F{0.11} F{0.27} F{0.16} F{0.24} Y}{\cab{ID} & \cab{Intercambio} & \cab{Dueño} & \cab{Contrato} & \cab{Falla}}
\mbox{INT-01} & Pedido preventa y consulta. & M3 & API negocio y sync. & Captura local sin confirmar. \\
\mbox{INT-02} & Entrega, POD y cobro. & M6/M7/M8 & API sync por lote. & Cola en dispositivo. \\
\mbox{INT-03} & Eventos bodega a nube. & M2/M5 & Eventos por sitio. & Broker local 24 h. \\
\mbox{INT-04} & Detalle de cross-docking a la nube. & M2 & Evento de inventario. & Buffer y reconciliación. \\
\mbox{INT-05} & Eventos de temperatura. & M9 & Evento de excursión. & Bloqueo local. \\
\mbox{INT-12} & Cambios de datos a réplica. & M2 & CDC de bodega. & Desfase medido; reanudar. \\
\mbox{INT-13} & Identidad a sitio. & Seguridad & Políticas firmadas. & Asignaciones vigentes. \\
\mbox{INT-14} & Métricas y trazas. & Operación & OTLP. & Buffer local. \\
\end{tablalafrox}

{\footnotesize Fuente: elaboración propia a partir de RT-02.06 y RT-03.12 (PUCV, 2026b, cap.~2, p.~7, y cap.~3, pp.~8--9).\par}

% Fichas del Anexo G. Compilación mediante anexos_4.1.tex desde la raíz.
\subsection*{Parámetros comunes del catálogo}
Los contratos distinguirán aceptación durable de procesamiento concluido. Como valores iniciales de diseño, las consultas interactivas tendrán espera máxima de 5 segundos y las transferencias de lotes, 30 segundos; un timeout produce estado incierto, nunca aprobación. Los reintentos conservan UUID, aplican espera exponencial con dispersión y máximo de 5 minutos entre intentos. Los errores de validación pasan a excepción sin reintento automático. El agotamiento del presupuesto de intentos deriva a una cola de errores con alerta; no elimina la operación. Estos parámetros se ajustarán con la prueba de carga sin relajar los tiempos de negocio exigidos.

\paragraph{INT-01. Preventa y consultas --- M3}
Modo mixto: consultas síncronas y captura/sincronización asíncrona. Volumen de pedidos: 1.400/2.600 por día; las consultas de lectura se estiman en cuatro por pedido (5.600/10.400 llamadas/día) y se miden por separado. Contraparte: API central, ventana requerida 24$\times$7; la app conserva 14 horas de trabajo. Dentro de 2 segundos, si no hay respuesta vigente, muestra la copia fechada; no confirma reserva ni crédito. Lotes de sincronización esperan hasta 30 segundos y se reenvían con la misma clave. La confirmación final corresponde a M2/M3 y al control de crédito.

\paragraph{INT-02. Reparto, POD y rendición --- M6/M7/M8}
Modo asíncrono para evidencia y rendición; consulta de resultado síncrona. Base mínima: 1.400/2.600 entregas diarias (caso, cap.~14). Se consideran cinco operaciones por entrega: 7.000/13.000 operaciones/día; los eventos no aplicables no se emiten. Fotografías y firmas viajan como objetos con hash, separados del mensaje. Contraparte: API y almacenamiento central, requeridos 24$\times$7. Timeout de lote: 30 segundos. La cola cifrada del dispositivo conserva el turno de 14 horas; solo elimina el evento local después de confirmar mensaje y objetos completos. No registra pago autorizado por ausencia de respuesta del POS.

\paragraph{INT-03. Eventos de bodega a nube --- M1/M2/M5/M9}
Modo asíncrono mediante outbox, RabbitMQ y consumidor canónico. Sobre de diseño: 260.000 ÷ 22,14 $\times$ 5 $\approx$ 58.710 eventos/día; peak: 109.032, con factor 1,857. Es el total de trazabilidad de los sitios, no una carga adicional por sitio. Contraparte: ingesta central, requerida 24$\times$7. Sin acuse en 30 segundos, el productor retiene y reintenta; el acuse del broker no sustituye el acuse transaccional del dominio. Se conserva al menos 24 horas de actividad local y se mide la edad del evento local más antiguo. El consumidor deduplica por sitio y UUID antes de actualizar el dominio.

\paragraph{INT-04. Detalle de cross-docking a la nube --- M2}
Modo asíncrono. Volumen de diseño: 5.600/10.400 eventos/día, derivado de cuatro operaciones por cada una de las 1.400/2.600 entregas; no se suma de nuevo al total de INT-03. Contraparte: consolidación N-04/N-05 en nube, requerida 24$\times$7. Timeout de acuse de 30 segundos; cada sitio mantiene su operación, buffer de 24 horas y secuencia propia. Al reconectar se reconcilia inventario en tránsito y se aísla cualquier doble compromiso.

\paragraph{INT-05. Temperatura --- M9}
Modo asíncrono para muestras y alarma local inmediata. Volumen: 10.584 lecturas/día en régimen y peak, por 21 puntos de cámara × 288 lecturas más 28 termógrafos × 162 lecturas entre 05:30 y 19:00 (6.048 + 4.536). Contraparte: IoT Core/ingesta central, requerida 24$\times$7; el bloqueo M9--M5 permanece local y no espera a nube. Acuse remoto máximo de 30 segundos por lote. Gateway: retención de 24 horas; termógrafo: registro de toda la ruta, hasta 14 horas. Una pérdida de señal o lectura obsoleta genera una incidencia diferenciada de la excursión térmica.

\paragraph{INT-12. Réplica de lectura del WMS de Talca --- M2/Datos}
Modo asíncrono CDC. El volumen físico depende del WAL y no equivale al número de eventos EPCIS. Para la prueba se usa $V_{WAL}=N_T\times b_{WAL}$, donde $N_T$ son los cambios medidos en Talca y $b_{WAL}$, bytes medidos por cambio. Los 12.602/23.404 cambios diarios son una cota conservadora calculada con los movimientos de todos los sitios, no un conteo propio de Talca; se aplica como cota al WAL de Talca, estimado en 0,123 GB/día. Migraciones, índices y amplificación se miden aparte. Contraparte: réplica Aurora de solo lectura, requerida 24$\times$7. Se alerta al superar 5 minutos de retraso y se escala antes de 15; si caen fibra y LTE, Starlink en espera caliente toma el tráfico prioritario de DMS/WAL y sostiene RPO $\leq$ 15 minutos; con los tres caminos caídos, el WAL se conserva localmente para reconciliación. AL-DR-01 verifica el escenario de pérdida de sitio.

\paragraph{INT-13. Identidad y manifiestos --- Seguridad}
Modo mixto: autenticación central síncrona y distribución asíncrona de manifiestos. Se estiman 760 eventos/día (380 dispositivos × 2), tanto en régimen como en peak. Contraparte: Keycloak, requerida 24$\times$7. Timeout interactivo: 5 segundos. Sin enlace se usa el verificador local con manifiesto de 26 horas, caché de identidad de 24 horas y permisos de turno de 8/14 horas. No se habilitan nuevas altas ni privilegios. AL-OFF-01 prueba dos relevos y el rechazo de identidades no preinscritas.

\paragraph{INT-14. Métricas, logs y trazas --- Operación}
Modo asíncrono por OTLP. Presupuesto de diseño: 13.000 eventos/día en régimen y peak (13 nodos × 1.000 eventos: 6 VM de Talca, 4 de Concepción y 3 mini-PC); son unidades de observabilidad y no mensajes comerciales sumables. Contraparte: colectores y CloudWatch, requeridos 24$\times$7. Exportación con timeout de 30 segundos y buffer en disco de 24 horas. El fallo de exportación no bloquea negocio; el sitio mantiene alarma local por ocupación. Auditoría de negocio y eventos de seguridad no se descartan para conservar trazas diagnósticas.

{\footnotesize Fuente: elaboración propia a partir de PUCV (2026a, cap.~14, p.~24) y RT-05.21 (PUCV, 2026b, cap.~5, p.~13). Las tasas de consultas son supuestos de ensayo explícitos; su calibración exige medición.\par}
